\documentclass[10pt,a4paper]{amsart}
\usepackage[margin=19mm]{geometry}
\usepackage{amsmath,amssymb,mathtools}
\usepackage[scr=rsfs,cal=euler]{mathalfa}
\usepackage{xfrac}
\usepackage[unicode,hidelinks,bookmarks=false]{hyperref}
\newcommand{\ie}{i.e.\ }
\newcommand{\cf}{cf.\ }
\newcommand{\resp}{respectively\ }
\newcommand{\Subsection}{\subsection}
\providecommand{\nobreakdash}{\nobreak\hbox{-}}
\setlength{\emergencystretch}{2em}
\multlinegap0pt
\usepackage{amssymb}
\usepackage{xcolor}
\usepackage[shortlabels]{enumitem}
\usepackage{float}
\newtheorem{assumption}{Assumption}
\newtheorem{theorem}{Theorem}
\newtheorem{lemma}{Lemma}
\title{A Statistical Framework \\ for Learning Preferences from the Past}
\author{Tamojit Sadhukhan, Moulinath Banerjee, Krishanu Maulik, Parthanil Roy}
\date{Source: arXiv:2605.10042v1 (CC BY 4.0)}
\begin{document}
\maketitle

\noindent\emph{Source and licence.} A Statistical Framework \\ for Learning Preferences from the Past. Version arXiv:2605.10042v1 (CC BY 4.0), distributed by arXiv under the Creative Commons Attribution 4.0 International licence, http://creativecommons.org/licenses/by/4.0/, as recorded in the arXiv licence metadata for this exact version and shown on the abstract page https://arxiv.org/abs/2605.10042v1. Full licence notice: \texttt{licenses/arxiv-2605.10042-CC-BY-4.0.txt}.\par\smallskip

\noindent\textit{Editorial scope.} Counted unit: the article's lemma lem1 (source0.tex lines 1136--1138). The article's complete original proof of it is delivered with it (source0.tex lines 1140--1150, one \texttt{proof} environment of 11 lines). No mathematics is written, repaired or re-derived: the statement and the proof are the article's own lines, in the article's own notation, and no retained line is substituted or re-flowed.  Retained uncounted as the source-local context the proof needs: lines 238--240; lines 254--256; lines 312--319.  Nothing was omitted from any retained span, so the package carries no omission record.  No retained line contains a citation, so the wrapper carries no bibliography and no bibliography entry.  No retained line contains a figure environment, an image-inclusion command or drawing code, so no image is delivered and the package declares no asset dependency.  Editorial formatting only, in addition: an AMS \texttt{amsart} preamble that restates the article's own declarations and macros used by the retained lines, namely the environments \texttt{assumption}, \texttt{theorem}, \texttt{lemma} on separate counters, together with the standard packages those lines need.  The retained environments print in this document as Assumption 1, Theorem 1, Lemma 1 in order of appearance, whereas the article prints the same items with the numbering of its own full document.  The complete native source of the article as distributed in the arXiv e-print for this exact version is bundled unchanged as \texttt{sources/200276/source0.tex}.
\begin{assumption}\label{a1}
   The positive integer valued random variable $T$ has bounded support. In other words, there exists a $T_0 \in \mathbb{N}$ such that $1 \leq T \leq T_0$ almost surely.
\end{assumption}

\begin{assumption}\label{a4}
    The function $h$ is continuously differentiable in $\mathcal{N}(r_0)$ with $h'(r_0) \neq 0$.
\end{assumption}

\begin{theorem}\label{thm2} Under Assumptions~\ref{a1} - \ref{a4}, we have
\begin{align*}
    \left(\Omega_N, \Omega^{(0)}_N\right) {\to} \left(g_{\alpha,\beta}, g^{(0)}_{\alpha,\beta}\right), \quad \text{ where } \quad 
    \alpha^2 = \frac{h\left(r_{0}\right)\left(1-h\left(r_{0}\right)\right)}{\mathcal{F}^\prime\left(r_{0}\right)}, \quad \beta = \frac{1}{2}h'(r_0),
\end{align*}
where the convergence is 
in the space $\mathcal{L}_2 \times \mathcal{L}_2$.
\end{theorem}

\begin{lemma}\label{lem1}
    For any $\epsilon > 0$, there exists $L_\epsilon > 0$ such that with probability at least $1-\epsilon$, $\widetilde{E}_{N} \subseteq [-L_\epsilon,L_\epsilon]$ for all sufficiently large $N$.
\end{lemma}

\begin{proof}
By definition, $\widetilde{E}_{N}$ is either the null set, or an interval $[X_N, Y_N]$ containing the origin, and in the first case the proof is immediate. For the latter case, fix $\epsilon > 0$. 
For any $M > 0$, $Y_N > M$ implies that $r_0 + N^{-1/3}M \in E_N$. Note that $r \in E_N$ implies either i) $\hat{h}_N(r) = \hat{h}_N(r_0)$ or ii) either $\hat{h}^{(0)}_N(r) = h_0 < \hat{h}_N(r) < \hat{h}_N(r_0)$ or $\hat{h}_N(r_0) < \hat{h}_N(r) < h_0 = \hat{h}^{(0)}_N(r)$. So,
\begin{equation*}
\{Y_N > M\} \subset  \{\Omega_N(M) = \Omega_N(0)\} \cup  \{\Omega^{(0)}_N(M) = 0\}.
\end{equation*}
By Theorem \ref{thm2}, for all sufficiently large $M$, we have
\begin{align*}&\mathbb{P}(\Omega_N(M) = \Omega_N(0)) \to \mathbb{P}(g_{\alpha,\beta}(M) = g_{\alpha,\beta}(0)) < \epsilon/4, \\
&\mathbb{P}(\Omega^{(0)}_N(M) = 0) \to \mathbb{P}(g^{(0)}_{\alpha,\beta}(M) = 0) = \mathbb{P}(g^+_{\alpha,\beta}(M) \leq 0)  < \epsilon/4.\end{align*}
Thus there exists $L_\epsilon > 0$ such that $\mathbb{P}(Y_N > L_{\epsilon}) < \epsilon/2$ for all sufficiently large $N$. In a similar way, we can show (by choosing $L_\epsilon$ so that it works in both cases) $\mathbb{P}(X_N < -L_{\epsilon}) < \epsilon/2$ for all sufficiently large $N$. Thus, $\mathbb{P}(\widetilde{E}_{N} \subseteq [-L_\epsilon,L_\epsilon]) > 1-\epsilon$ for all sufficiently large $N$.
\end{proof}

\end{document}
